\SourceLeaf{426}%
die Fläche der Wert von $V$ sein Vorzeichen ändern, sodass, solange die Stetigkeit gewahrt bleibt, auf der einen Seite die positiven und auf der anderen die negativen [Werte] auftreten.

Die Richtung der Normalen zur Fläche zu der Seite hin, die wir als obere betrachten, werde auf der Hilfskugel durch den Punkt $L$ bezeichnet, und es gelte
\[
\cos L(1)=X,\qquad\cos L(2)=Y,\qquad\cos L(3)=Z.
\]
Ferner sei $ds$ eine unendlich kleine Linie auf der Fläche, deren Richtung auf der Hilfskugel durch den Punkt $\lambda$ bezeichnet wird, und es gelte
\[
\cos\lambda(1)=\xi,\qquad\cos\lambda(2)=\eta,\qquad\cos\lambda(3)=\zeta.
\]
Dann haben wir
\[
dx=\xi ds,\qquad dy=\eta ds,\qquad dz=\zeta ds,
\]
also
\[
P\xi+Q\eta+R\zeta=0
\]
und, da $\lambda L=90^{0}$ gelten muss, auch
\[
X\xi+Y\eta+Z\zeta=0.
\]
Da $P,Q,R,X,Y,Z$ nur von der Stelle der Fläche abhängen, von der das Element ausgeht, und diese Gleichungen für jede Richtung des Elements in der Fläche gelten, sieht man leicht, dass $P,Q,R$ zu $X,Y,Z$ proportional sein müssen, also
\[
P=X\mu,\qquad Q=Y\mu,\qquad R=Z\mu,
\]
Da nun
\[
XX+YY+ZZ=1
\]
gilt, erhält man
\[
\mu=PX+QY+RZ
\]
und
\[
\mu\mu=PP+QQ+RR
\]
oder
\[
\mu=\pm\BareRoot{(PP+QQ+RR)}.
\]

Entfernt man sich von der Fläche in Richtung der Normalen um\space%
